\chapter{``朴素''作用量中的费米子加倍问题}

由式~7.7 中的算符 $M^N$ 所刻画的狄拉克作用量的朴素离散化，其问题在于：在连续极限下它给出 $2^d = 16$ 味而非一味。这一加倍问题很容易通过自由场传播子之逆（对所有 $U_\mu(x) = 1$ 的作用量作傅里叶变换得到）来演示。
\begin{equation}
S^{-1}(p) = m_q + \frac{i}{a}\sum_\mu \gamma_\mu \sin p_\mu a
\end{equation}
在 $m_q \to 0$ 极限下，该式在布里渊区内有 16 个零点。若将布里渊区的动量范围定义为 $\{-\pi/2, 3\pi/2\}$，则这些零点位于 $p_\mu = 0$ 和 $\pi$ 处。正如后文所讨论的，这种增殖在 Nielsen-Ninomiya 定理 [14] 所规定的非常一般的条件下成立，并且与同时保持手征对称性密切相关。包含规范场并不能解决加倍问题。

在讨论部分修正之前，考察这些额外的零模在手征变换下的性质是有益的。让我们定义一组 16 个四矢量 $\Pi^A = \{(0, 0, 0, 0), \allowbreak (\pi, 0, 0, 0), \allowbreak \ldots, \allowbreak (\pi, \pi, \pi, \pi)\}$，其中 $A = \{1 \ldots 16\}$，并考虑无质量传播子在这些点附近的展开。于是
\begin{equation}
\begin{aligned}
S^{-1}(p, m = 0) &= \frac{i}{a}\sum_\mu \gamma_\mu \sin p_\mu a \\
&= \frac{i}{a}\sum_\mu \gamma_\mu \sin(\Pi^A + k)_\mu a \\
&= \frac{i}{a}\sum_\mu \gamma_\mu S_\mu^A \sin k_\mu a \\
&\equiv \frac{i}{a}\sum_\mu \tilde{\gamma}_\mu \sin k_\mu a ,
\end{aligned}
\end{equation}
其中 $S_\mu^A = \{+1, -1\}$，取决于 $\Pi^A$ 的 $\mu$ 分量是 $\{0, \pi\}$ 中的哪一个。在最后一个表达式中，通过相似变换 $\tilde{\gamma}_\mu = Y_A \gamma_\mu Y_A^\dagger$ 定义了 $\gamma$ 矩阵的一个新表示，其中
\begin{equation}
Y_A = \prod_\mu (\gamma_\mu \gamma_5)^{n_\mu^A}
\end{equation}
而 $n_\mu^A$ 取 $\{0, 1\}$，取决于 $\mu$ 方向上的动量展开是绕 0 还是绕 $\pi$。于是，$\tilde{\gamma}_5 = Y_A \gamma_5 Y_A^\dagger =
S_1^A S_2^A S_3^A S_4^A \gamma_5 \equiv X^A \gamma_5$，其中符号 $X^A$ 列于表 3。因此，十六个种态按手征荷 $\pm 1$ 分裂为两组各 8 个，从而使理论成为无反常的。因此，这种``朴素''离散化在现象学上是不可接受的。事实上，加倍子的存在可以追溯到朴素费米子作用量的一个更大的味对称性，下面我来讨论这一点。

\begin{table}[htbp]
\centering
\begin{tabular}{ccc}
\toprule
$A$ & 简并度 & $X^A$ \\
\midrule
$(0, 0, 0, 0)$ & 1 & $+1$ \\
$(1, 0, 0, 0)$ & 4 & $-1$ \\
$(1, 1, 0, 0)$ & 6 & $+1$ \\
$(1, 1, 1, 0)$ & 4 & $-1$ \\
$(1, 1, 1, 1)$ & 1 & $+1$ \\
\bottomrule
\end{tabular}
\caption{加倍子的手征符号与简并因子。}
\label{tab:3}
\end{table}

在连续情形下，无质量狄拉克作用量在 $U(1)_V \otimes U(1)_A$ 味变换下不变
\begin{equation}
\begin{aligned}
\psi_L &\to V_L \psi_L , & \bar{\psi}_L &\to \bar{\psi}_L V_L^\dagger , \\
\psi_R &\to V_R \psi_R , & \bar{\psi}_R &\to \bar{\psi}_R V_R^\dagger .
\end{aligned}
\end{equation}
由于左手场与右手场仅通过质量项耦合，矢量对称性 $V_L = V_R = e^{i\theta}$ 对所有 $m$ 都成立，而轴矢对称性 $V_L = V_R^\dagger = e^{i\theta}$（或等价地 $V_L = V_R = e^{i\theta\gamma_5}$）仅对 $m = 0$ 成立。单味的朴素费米子作用量 $S_N$ 在下面变换下具有大得多的对称群 $U(4)_V \otimes U(4)_A$：
\begin{equation}
\begin{aligned}
\psi(x) &\to \Gamma_x \left[ \frac{1 - \epsilon(x)\gamma_5}{2} V_L + \frac{1 + \epsilon(x)\gamma_5}{2} V_R \right] \Gamma_x^\dagger \psi(x) \\
\bar{\psi}(x) &\to \bar{\psi}(x) \Gamma_x \left[ V_L^\dagger \frac{1 + \epsilon(x)\gamma_5}{2} + V_R^\dagger \frac{1 - \epsilon(x)\gamma_5}{2} \right] \Gamma_x^\dagger
\end{aligned}
\end{equation}
前提是考虑间隔 $2a$ 的平移。$\Gamma_x$ 由式~11.1 给出，相位 $\epsilon(x)$ 由式~11.4 给出，并且 $\Gamma_x \gamma_5 \Gamma_x^\dagger = \epsilon(x)\gamma_5$。这里 $V_{L,R}$ 是作用在费米子旋量指标上的一般 $4 \times 4$ 幺正矩阵。尽管在式~9.5 中通过 $\epsilon(x)$ 存在对 $x$ 的依赖，但这是一个整体对称性，因为在一个点上将其给定便在所有点上将其确定。

为了看清加倍是如何产生的，考虑矢量（$V_L = V_R$）离散子群，其三十二个元素 $\{\pm Y_A\}$ 由式~9.3 定义。式~9.2 表明，$Y_A$ 沿 $A$ 所指定的方向将动量平移 $\pi/a$，
\begin{equation}
Y_A S_F(k) Y_A^\dagger = S_F\!\left( k + \frac{\pi}{a} A \right) .
\end{equation}
由于这是一个相似变换，由该变换联系起来的布里渊区的十六个区域在物理上是等价的。这一分析仅依赖于作用量的旋量结构，使得我们容易理解为什么即使在相互作用理论中加倍子依然存在。在第 11 节中，这一对称性将被利用来构造交错费米子。在该形式中，十六个加倍子被约化为四个，同时保留了一个 $U(1)_V \otimes U(1)_\epsilon$ 对称性（对应于 $V_{L,R}$ 只是相位的情形）。$U(1)_\epsilon$ 在有限的 $a$ 下扮演手征对称性的角色，并在连续极限下变为 $U(1)_A$。

\section{手征对称性的一般性质}

在连续理论的无质量极限下，经典层面的味对称性为 $U_V(1) \times U_A(1) \times SU_V(n_f) \times SU_A(n_f)$，其中 $n_f$ 是轻费米子的数目。未破缺的 $U_V(1)$ 给出重子数守恒。$U_A(1)$ 被瞬子贡献所破缺，且味单态轴矢流是反常的。$SU_A(n_f)$ 发生自发破缺，夸克凝聚在 QCD 真空中获得非零期望值。与这一自发破缺对称性相联系的是一批 $(n_f^2 - 1)$ 个 Goldstone 玻色子，即$\pi$介子。手征对称性在连续理论中扮演了非常有用的角色。例如，人们可以（i）在 $SU_L(n_f) \times SU_R(n_f)$ 的不同表示下对算符进行分类，而手征对称性禁止不同表示中的算符之间发生混合；（ii）推导给定算符在末态中$\pi$介子数目不同的各态之间的矩阵元关系（例如 $K \to \pi\pi \leftrightarrow K \to \pi\pi\pi$）；（iii）利用算符沃德恒等式；以及（iv）推导关联函数之间的关系，即沃德-高桥恒等式。因此，在有限的 $a$ 下（而不仅仅是在连续极限下）保持手征对称性是人们所期望的。要确定这是否可能，就需要理解格点理论中加倍子与手征对称性之间的关系。

式~7.11 所定义的手征对称性由格点理论实现，只要 $\gamma_5 M + M \gamma_5 = 0$。容易验证，朴素格点作用量满足厄米性性质
\begin{equation}
\gamma_5 M \gamma_5 = M^\dagger
\end{equation}
因此，在无质量极限下，欧几里得格点作用量必须是反厄米的，手征对称性（式~7.11）才能成立。式~7.7 中给出的 $M^N$ 正是如此，而这只需对导数取对称差分即可实现。重要的问题是——在什么一般条件下，具有这种手征对称性实现的格点理论会出现加倍子。答案由 Nielsen-Ninomiya 定理 [14] 给出。考虑一个推广的作用量，使得 $S^{-1}(p, m = 0) = iF(p)$。于是，若
\begin{itemize}
\item 函数 $F(p)$ 在动量空间中具有周期 $2\pi/a$。（这是平移不变性的一个结果。）
\item 格点动量在 $\{0, 2\pi\}$ 范围内连续。这在所有格距取值下于 $L \to \infty$ 极限中成立。
\item $F(p)$ 在动量空间中连续。若相互作用是定域的，这一点便得到保证。
\item 对于小的 $p_\mu$ 且当 $a \to 0$ 时，$F(p) \to p_\mu \gamma_\mu$，以与连续理论匹配。
\item 作用量具有手征对称性，即它满足厄米性性质 $\gamma_5 M \gamma_5 = M^\dagger$，且无质量狄拉克作用量是反厄米的。
\end{itemize}
则该理论将具有加倍子。这一点很容易从图~6 所示的 $F(p)$ 的示意性一维图中理解，该图突出了在上述条件下 $F(p)$ 必须具有偶数个零点这一事实。

\begin{figure}[htbp]
\centering
\includegraphics[width=0.8\textwidth]{images/fig6.png}
\caption{无质量逆传播子 $F(p)$ 的示意性一维表示。}
\label{fig:6}
\end{figure}

$\pm$ 斜率意味着一个态向左运动，而另一个态向右运动。在一维情形下，运动方向与手征性相同，因此左手场和右手场以不可分离的对出现，并共同构成一个狄拉克费米子。

Nielsen-Ninomiya 定理的关键结论是：人们无法定义一个具有手征费米子的格点理论。例如，若为一个左手狄拉克费米子写下 $M^N$，那么它的加倍子将是右手的，但在 $L \otimes R$ 下具有相同的变换。因此，这两个``态''共存而构成一个狄拉克费米子。假定 Nielsen-Ninomiya 定理是严密的，那么人们如何继续推进（i）定义一个没有加倍子并在连续极限下恢复手征对称性的格点理论，以及（ii）在格点上定义手征规范理论？

第一个问题的答案将通过显式重建 Wilson 修正与交错修正给出。另一个非常令人振奋的新进展——我无法在此充分阐述——是用于手征对称表述的 Ginsparg-Wilson 条件的复活 [53]，即 [54]
\begin{equation}
\gamma_5 D + D\gamma_5 = a D \gamma_5 D .
\end{equation}
L\"uscher 最近给出了一份清晰的分析，说明这一条件如何允许构造一个没有加倍子且具有手征对称性的格点理论，如何避开 NN 定理，并给出 ABJ 反常 [55]。遗憾的是，这一替代方案的瓶颈在于缺少足够定域从而使模拟可行的格点狄拉克作用量的闭式表达式。第二个问题——构造手征规范理论——仍未解决。对于那些对这一非常重要的未解问题感兴趣的人，评述文章 [15--17] 是一个很好的起点，其中包含了一些最近的尝试。
